\chapter{Strep throat}
\index{Strep throat}

\medskip
\noindent\textit{Educational reference; always seek professional advice for individual care.}

\section*{Definition and Overview}
Strep throat is an infection of the throat and tonsils caused by the bacterium \textit{Streptococcus pyogenes}, also called group A streptococcus. It is one of the commonest causes of a sore throat in school-aged children and adolescents, though it can affect people of any age. Sore throats are extremely common and most are caused by viruses, which do not respond to antibiotics, so the importance of strep throat lies in telling it apart from viral infections and in treating it where appropriate. Antibiotic treatment of strep throat shortens the illness, reduces the chance of spreading it to others, and, most importantly, prevents rare but serious complications such as rheumatic fever and kidney disease. In many countries, rheumatic fever remains a significant problem in Aboriginal and Torres Strait Islander communities and in parts of the Northern Territory, Queensland, and western regions, which is one reason doctors in those regions treat suspected strep throat more readily. By contrast, in most of the community, antibiotics are not needed for every sore throat, and unnecessary use drives antibiotic resistance and side effects. This chapter explains the features of strep throat, how it is diagnosed, when antibiotics are warranted, how it is treated at home, and the warning signs that require urgent care.

\section*{Epidemiology}
Sore throat is one of the most frequent reasons for medical visits, and strep throat accounts for a significant minority of cases: roughly 5 to 15 percent of sore throats in adults and 20 to 30 percent in children are caused by group A streptococcus. The infection is most common between the ages of five and fifteen, is uncommon under the age of three, and occurs slightly more often in winter and early spring in temperate some countries. It spreads easily through respiratory droplets produced by coughing and sneezing, and through shared cups, utensils, and hands, so it moves quickly through schools, childcare centres, and households. Asymptomatic carriers, people who carry the bacterium without symptoms, are common in childhood and complicate the interpretation of throat swabs. Rheumatic heart disease, the most feared long-term complication, is now rare in wealthy countries, but the affected region retains a disproportionately high burden among Aboriginal and Torres Strait Islander people, which is why national guidelines recommend a low threshold for treating suspected streptococcal throat infections in high-risk communities.

\section*{Aetiology and Pathophysiology}
Group A streptococcus, the bacterium responsible for strep throat, is carried in the nose and throat of healthy people and transmitted by droplets or direct contact. Once inhaled or transferred to the throat, the bacteria attach to the lining of the tonsils and pharynx and release toxins and enzymes that damage tissue, producing the classic red, swollen, and sometimes pus-covered throat.

\begin{itemize}
    \item \textbf{\textit{Streptococcus pyogenes}:} the sole bacterial cause of true strep throat; a gram-positive coccus that grows in chains
    \item \textbf{Direct spread:} respiratory droplets from coughing, sneezing, or talking, and contact with contaminated hands or objects
    \item \textbf{Carriage:} some people, especially children, carry the bacteria without illness and can pass it on
    \item \textbf{Bacterial toxins:} released enzymes damage the throat lining and drive fever and inflammation; certain toxin-producing strains cause the rash of scarlet fever
    \item \textbf{Immune complications:} the body's antibody response to the bacteria can, in a minority, attack the heart valves (rheumatic fever) or the kidneys (post-streptococcal glomerulonephritis)
    \item \textbf{Environmental factors:} crowding, sharing drinking vessels, and close contact in schools and families increase transmission
\end{itemize}

\section*{Clinical Features}
Strep throat typically begins suddenly, and certain features distinguish it from the common viral sore throat.

\begin{itemize}
    \item \textbf{Sudden severe sore throat:} pain that is worse on swallowing and often intense from the start
    \item \textbf{Fever:} temperature of 38\textdegree{}C or higher, sometimes with chills
    \item \textbf{Red, swollen tonsils:} often covered with white patches or streaks of pus
    \item \textbf{Swollen neck glands:} tender lymph nodes at the angles of the jaw
    \item \textbf{Absence of cough, runny nose, and hoarseness:} these viral features make strep throat less likely
    \item \textbf{Headache and general malaise:} feeling generally unwell, tired, and achy
    \item \textbf{Scarlet fever rash:} a fine, red, sandpaper-like rash on the trunk and limbs with a strawberry tongue, occurring in some children
    \item \textbf{Nausea, vomiting, or abdominal pain:} more common in younger children
    \item \textbf{Petechiae:} tiny red spots on the roof of the mouth, a characteristic but not always present sign
\end{itemize}

\section*{Differential Diagnosis}
A sore throat has many causes, and the diagnosis of strep throat rests on distinguishing it from these.

\begin{itemize}
    \item \textbf{Viral pharyngitis:} the commonest cause; typically accompanied by cough, runny nose, and hoarseness, with a more gradual onset and no pus on the tonsils
    \item \textbf{Infectious mononucleosis:} glandular fever caused by Epstein-Barr virus; severe sore throat, fatigue, and swollen glands, common in teenagers and young adults
    \item \textbf{Tonsillitis from other bacteria:} \textit{Fusobacterium} (the cause of quinsy), and rarely other streptococcal species, can look similar
    \item \textbf{Common cold and influenza:} viral illnesses with prominent cough and congestion rather than severe throat pain
    \item \textbf{Gastro-oesophageal reflux:} throat irritation and pain after eating or lying down, without fever or swollen glands
    \item \textbf{Oral thrush or mouth ulcers:} visible lesions without the classic pattern of tonsillar pus and tender neck glands
\end{itemize}

\section*{When to Seek Medical Care}
Most sore throats settle within a few days without medical treatment, but a doctor should assess a sore throat that is severe, lasts more than a few days, or is accompanied by high fever, difficulty swallowing fluids, or a rash. Medical review is particularly important for children in high-risk communities where rheumatic fever prevention is a priority, for people with weakened immune systems, and for anyone whose symptoms are worsening. A doctor can perform a swab and, if appropriate, prescribe antibiotics. There is no need to see a doctor for every mild sore throat, but young children who cannot take fluids, and anyone who feels unusually unwell, should be reviewed. Parents and carers can trust their instincts: a child who is normally lively but becomes quiet, refuses drinks, and has a high fever deserves prompt assessment even if the throat itself does not look severe.

\textbf{Contact your local emergency services for:}
\begin{itemize}
    \item Difficulty breathing, or noisy, raspy breathing
    \item Severe difficulty swallowing, including drooling and an inability to swallow saliva
    \item Muffled voice with a "hot potato" sound, which can indicate a throat abscess
    \item A stiff neck, severe headache, sensitivity to light, or a seizure, which can indicate meningitis
    \item Extreme lethargy, confusion, or a rapid, spreading rash
\end{itemize}

\section*{Diagnosis and Investigations}
The diagnosis of strep throat combines clinical assessment with, when needed, laboratory confirmation.

\begin{itemize}
    \item \textbf{Clinical scoring:} features such as fever, tonsillar pus, swollen neck glands, and absence of cough are scored; patients with several features are more likely to have strep throat
    \item \textbf{Throat swab and culture:} a swab rubbed over the tonsils is cultured in the laboratory; it is the reference standard but takes a day or two
    \item \textbf{Rapid antigen test:} detects group A streptococcus from a throat swab within minutes; used in some clinics to guide decisions
    \item \textbf{Blood tests:} sometimes used to check for glandular fever, or to assess complications such as kidney involvement
    \item \textbf{Routine swabs in high-risk regions:} in areas with high rates of rheumatic fever, swabbing and treating is performed more liberally under local guidelines
\end{itemize}

\section*{Management}
Treatment aims to relieve symptoms, reduce spread, and prevent complications.

\begin{itemize}
    \item \textbf{Antibiotics:} phenoxymethylpenicillin (or amoxicillin) is the standard treatment, given for about ten days to fully clear the bacteria; azithromycin is used in penicillin allergy
    \item \textbf{Pain relief:} paracetamol or ibuprofen reduces fever and throat pain; follow dose instructions carefully, especially for children
    \item \textbf{Fluids and rest:} cool drinks and soft foods such as soup and ice-cream are soothing, and rest supports recovery
    \item \textbf{Throat lozenges and gargles:} lozenges, or salt-water gargles, can ease pain in older children and adults, but honey and lozenges are not recommended for children under one
    \item \textbf{Analgesic mouth sprays:} short-term numbing sprays can help severe pain in adults
    \item \textbf{Staying home:} children and adults with confirmed strep throat should stay home until they have been on antibiotics for 24 hours or are feeling better, to reduce spread
    \item \textbf{Follow-up:} recheck is advised if symptoms do not improve within 48 hours or return after treatment
    \item \textbf{Treating complications:} an abscess may need drainage by a specialist, and rheumatic fever requires long-term preventive penicillin
\end{itemize}

\section*{Complications}
Complications of strep throat are uncommon but can be serious, and are the main reason for antibiotic treatment.

\begin{itemize}
    \item \textbf{Peritonsillar abscess:} a collection of pus behind the tonsil causing severe pain, trismus (difficulty opening the mouth), and muffled voice; needs drainage and intravenous antibiotics
    \item \textbf{Acute rheumatic fever:} an immune reaction causing fever, arthritis, and inflammation that can permanently damage the heart valves, leading to rheumatic heart disease
    \item \textbf{Post-streptococcal glomerulonephritis:} kidney inflammation occurring weeks after infection, causing blood in the urine, swelling, and high blood pressure
    \item \textbf{Otitis media and sinusitis:} infection can spread to the middle ear and sinuses
    \item \textbf{Scarlet fever:} a toxin-related rash occurring with the throat infection; responds to the same antibiotic treatment
    \item \textbf{Invasive disease:} rarely, the bacteria enter the bloodstream, causing sepsis, toxic shock, or tissue destruction, which is a medical emergency
\end{itemize}

\section*{Prognosis and Outcomes}
The great majority of people with strep throat recover fully within a week, and symptoms usually begin to improve within 24 to 48 hours of starting antibiotics. Left untreated, the illness typically settles on its own over a week to ten days, but the risk of complications, especially rheumatic fever, is higher. With timely treatment, rheumatic fever is largely preventable, which is why national guidelines are careful about who receives antibiotics. In many countries's high-risk Aboriginal and Torres Strait Islander communities, public health programs emphasise early treatment of sore throats and long-term penicillin to prevent recurrence of rheumatic fever, and these measures have reduced the burden of rheumatic heart disease. For otherwise healthy people, strep throat rarely leaves any lasting problem, and repeat infections, though unpleasant, respond to the same approach. A small minority of people carry the bacteria after treatment and need a second course; this is harmless in the person themselves but can matter in households with high-risk members. During the recovery phase, good sleep, gentle activity, and adequate nutrition support the immune system and help prevent the lingering tiredness that can follow any significant infection.

\section*{Prevention and Self-Care}
\begin{itemize}
    \item Wash hands regularly with soap and water, and keep them away from the face
    \item Cover coughs and sneezes with a tissue or the elbow, and dispose of tissues promptly
    \item Avoid sharing cups, bottles, cutlery, and toothbrushes during an outbreak
    \item Keep children home from school or childcare while infectious and feeling unwell
    \item Take the full course of antibiotics exactly as prescribed, even when feeling better
    \item Use regular pain relief and plenty of fluids while recovering
    \item Seek medical review promptly in high-risk communities or for anyone with a weakened immune system
\end{itemize}

\section*{Sources and Further Reading}
The following organisations provide reliable information about sore throats, strep throat, and rheumatic fever.

\begin{itemize}
    \item NHS. \textit{Sore throat and tonsillitis: symptoms, treatment and when to see a doctor.} https://www.nhs.uk/
    \item Mayo Clinic. \textit{Strep throat: symptoms and causes.} https://www.mayoclinic.org/
    \item healthdirect Australia. \textit{Strep throat: symptoms and treatment.} https://www.healthdirect.gov.au/
    \item Australian Department of Health and Aged Care. \textit{Rheumatic fever and rheumatic heart disease strategy.} https://www.health.gov.au/
    \item World Health Organization (WHO). \textit{Rheumatic fever and rheumatic heart disease fact sheet.} https://www.who.int/
    \item RHDAustralia (Rheumatic Heart Disease Australia). \textit{ARF and RHD prevention guidelines.} https://www.rhdaustralia.org.au/
\end{itemize}

\clearpage
